\section{Related Work}
\label{sec:related}

\para{Prompt format sensitivity.} Early work on in-context learning showed that accuracy depends on the choice and order of demonstrations~\cite{zhao2021calibrate,lu2022fantastically}, and that models can do as well with irrelevant or misleading instructions as with good ones~\cite{webson2022prompt}. \citet{sclar2024quantifying} isolated \emph{formatting}: by jointly varying separators, spacing, and casing in few-shot completion prompts, they found spreads of tens of accuracy points that did not vanish with scale or instruction tuning. \citet{voronov2024mind} showed that template choice can outweigh the gains of proposed in-context learning methods. \citet{salinas2024butterfly} found that small changes, such as adding a space or changing the requested output format, flip a noticeable fraction of predictions in chat models, and \citet{he2024does} found sizeable differences between plain-text, Markdown, JSON, and YAML prompt templates for GPT models. Robustness benchmarks such as PromptBench~\cite{zhu2023promptbench} study adversarial character- and word-level perturbations. Unlike these studies, we apply one atomic edit at a time to a fixed chat-formatted template, which lets us attribute each effect to a specific edit. We also measure length and compliance in addition to accuracy, and we compare every effect with a measured noise floor.

\para{Multi-prompt evaluation.} Given this sensitivity, \citet{mizrahi2024state} argued that LLM benchmarks should report performance over a distribution of instruction paraphrases, and \citet{polo2024efficient} proposed estimating that distribution efficiently. Evaluation harnesses also stress that small template details affect reproducibility~\cite{biderman2024lessons}. Our results support this agenda and refine it: the need for multi-format evaluation is largest below about 1.5B parameters, while at 7--14B single-edit effects on classification are mostly near noise level.

\para{Instruction following and structured output.} \ifeval~\cite{zhou2023instruction} measures compliance with verifiable instructions such as word counts, required keywords, and letter case. We use its checkers to test whether casing changes compliance. Requesting structured output is another common practice. \citet{tam2024let} found that strict format constraints such as JSON can reduce reasoning accuracy. We separate two things that are often conflated: wrapping the \emph{input} in JSON and requesting JSON \emph{output}. We find that the second has much larger and task-dependent effects than the first.

\para{Nondeterminism in LLM inference.} Greedy decoding is often treated as deterministic, but floating-point results in batched inference depend on batch composition, so outputs can change between runs~\cite{atil2024llm,he2025defeating}. We quantify this directly for our stack (\vllm~\cite{kwon2023efficient}). Batch-dependent noise is negligible for short classification outputs but moves \ifeval compliance by up to about 3 points, which sets a lower limit on the effects we can detect for long generations.
